\begin{abstract}
Let \(\Phi\) be the classical Riemann, or de Bruijn--Newman, kernel. We prove
that every translation minor of \(\Phi\) of order at most four is strictly
positive at strictly ordered nodes, and an exact origin-confluent calculation
gives a negative order-five minor. Thus the exact global Polya-frequency order
is four. For the Riemann-kernel positive theorem, a sweep-free two-mode
argument proves global positivity of the log-curvature quantities \(q,F_2\)
and of the fully confluent invariant \(\Cfour\). Its finite certificate uses
only rational polynomial algebra, Taylor inequalities, and geometric tails.
The order-five origin certificate likewise uses exact rational arithmetic.
A
quotient-Wronskian argument reduces arbitrary fourth-order minors to
\(\partial_\xi\Psi<0\). In the increasing curvature coordinate, its numerator
is
\[
 \delta\Lambda\int_p^w
 W(t)\frac{\Cfour(t)}{Q(t)^6\kappa(t)^2}\dd t.
\]
Here \(W\) is the difference of two explicitly normalized triangular CDFs.
Its strict positivity follows from a one-crossing density ratio. The structural
transport reductions and analytic sign estimates are displayed in the text
and appendices; exact coefficient reconstructions remain explicitly identified
certificate boundaries.
We then construct an explicit positive even Schwartz kernel which is itself
strictly \(\PF_4\) but whose entire Fourier transform has nonreal zeros; the
the latter sign follows from a short exact discriminant inequality. Thus
continuous \(\PF_3\) and \(\PF_4\) membership alone do not force Fourier
real-rootedness. The results are finite-order and have no implication for the
Riemann Hypothesis.
\end{abstract}
